% Blueprint content for leslie2.
% This document builds up to the soundness of probabilistic forward simulation:
% if a system C is forward-simulated by a system A, then every trace distribution
% achievable by C is achievable by A.

\section{Systems, executions, and trace distributions}

Throughout, $Label$ is a fixed label type carrying a distinguished \emph{silent}
label $\tau$ (the type class \texttt{Silent}). A label is \emph{internal} when
it equals $\tau$ and \emph{external} otherwise; this partition depends only on the
label type, not on any individual system.

\begin{definition}[Probabilistic labelled transition system]
  \label{def:system}
  \lean{PLTS.System}\leanok
  A \emph{PLTS} over a state type $State$ and label type $Label$ is a designated
  initial state $init : State$ together with a step relation
  $step : State \to Label \to \operatorname{PMF} State \to \mathrm{Prop}$. We write
  $s \xrightarrow{l} \mu$ for $step\ s\ l\ \mu$: from $s$, on label $l$, the system
  may move to the probability distribution $\mu$ over successor states.
\end{definition}

\begin{definition}[Execution]
  \label{def:exec}
  \lean{PLTS.AlterSeq}\leanok
  An \emph{execution} is an alternating sequence $e = \langle init, trans\rangle$:
  an initial state followed by a possibly infinite sequence
  $trans : \operatorname{Seq}(Label \times State)$ of $(label, next\ state)$ pairs.
  Its \emph{trace} keeps only the external labels.
\end{definition}

\begin{definition}[Trace]
  \label{def:trace}
  \lean{PLTS.System.trace}\leanok
  \uses{def:exec}
  The trace of an execution $e$ is the sequence of labels obtained from $trans$ by
  deleting every internal ($\tau$) transition and projecting to the label
  component. Executions that differ only by trailing $\tau$-transitions have the
  same trace.
\end{definition}

\begin{definition}[Probabilistic execution and its path measure]
  \label{def:probexec}
  \lean{PLTS.ProbabilisticExecution}\leanok
  \uses{def:system, def:exec}
  A \emph{probabilistic execution} of $sys$ is an initial distribution
  $initState : \operatorname{PMF} State$ together with a \emph{scheduler} that, at
  each finite history, returns a sub-distribution over the next $(l, \mu)$ move (or
  halts). For a finite execution $e$ it induces a mass
  $probOf(e) \in [0, \infty]$
  (\lean{PLTS.ProbabilisticExecution.probOf}), the probability that the
  scheduler, started from $initState$, follows exactly the path $e$: the initial
  mass on $e.init$ times the product of the one-step kernels along $trans$.
\end{definition}

The trace probability collapses the path measure onto observable behaviour. To
avoid counting a single trajectory once for every trailing block of
$\tau$-transitions, we sum only over \emph{tight} executions.

\begin{definition}[Tight execution]
  \label{def:istight}
  \lean{PLTS.System.IsTight}\leanok
  \uses{def:exec}
  A finite execution is \emph{tight} if it has no transitions, or its last
  transition is external. Each trajectory has, for a given finite trace, a unique
  tight prefix: the shortest prefix with that trace (ending right after its last
  external transition).
\end{definition}

\begin{definition}[Trace probability]
  \label{def:traceprob}
  \lean{PLTS.System.traceProb}\leanok
  \uses{def:probexec, def:trace, def:istight}
  For a probabilistic execution $pe$ and a finite trace $\sigma$, the trace
  probability $traceProb(pe, \sigma)$ is the sum of $probOf(e)$ over all tight
  finite executions $e$ with trace $\sigma$. Tightness makes the summands
  correspond to disjoint cylinders, so this is the standard trace-cone
  probability: the chance that the run's external labels begin with $\sigma$.
\end{definition}

\begin{definition}[Achievable trace distributions]
  \label{def:achievable}
  \lean{PLTS.achievableTraceDists}\leanok
  \uses{def:traceprob}
  The set $\operatorname{ATD}(sys)$ of \emph{achievable trace distributions} of
  $sys$ is the set of functions $\sigma \mapsto traceProb(pe, \sigma)$ obtained
  from probabilistic executions $pe$ whose initial distribution is the Dirac
  $\delta_{sys.init}$ (genuine runs from the designated start state). Two systems
  are trace-distribution equivalent when these sets coincide.
\end{definition}

The Dirac-init restriction is essential: it is exactly what makes the
distribution-monad construction below trace-preserving, and it is the semantics in
which all of the soundness statements are phrased.

\section{Weak transitions and the two constructions}

Weak transitions let internal ($\tau$) computation be absorbed. All three are
distribution-to-distribution, run-to-halt notions.

\begin{definition}[Silent weak transition]
  \label{def:weaktau}
  \lean{PLTS.weakTau}\leanok
  \uses{def:system}
  $weakTau\ sys\ \mu_0\ \mu$ holds when there is an internal-only scheduler that,
  started from $\mu_0$, halts almost surely and whose distribution of final states
  is $\mu$. Intuitively $\mu_0 \Rightarrow_\tau \mu$: run $\tau$-steps until
  spontaneous termination, and $\mu$ records where you stop.
\end{definition}

\begin{definition}[Hyper-step]
  \label{def:hyperstep}
  \lean{PLTS.hyperStep}\leanok
  \uses{def:system}
  $hyperStep\ sys\ \mu_{\mathrm{pre}}\ l\ \mu_{\mathrm{post}}$ holds when
  $\mu_{\mathrm{post}}$ is a convex combination of one-step $l$-successors of the
  points of $\mu_{\mathrm{pre}}$: there is $p : State \to \operatorname{PMF}(\operatorname{PMF} State)$
  with $s \xrightarrow{l} \mu$ for every $s \in \operatorname{supp}\mu_{\mathrm{pre}}$
  and $\mu \in \operatorname{supp}(p\,s)$, and
  $\mu_{\mathrm{post}} = \mu_{\mathrm{pre}} \mathbin{\gg\!=} (s \mapsto (p\,s)\mathbin{\gg\!=}\mathrm{id})$.
\end{definition}

\begin{definition}[External weak transition]
  \label{def:weakstep}
  \lean{PLTS.weakStep}\leanok
  \uses{def:weaktau, def:hyperstep}
  $weakStep\ sys\ \mu_0\ l\ \mu_f$ is a $\tau^\ast\!\cdot l\cdot\tau^\ast$ chain:
  there exist $\mu, \mu'$ with $\mu_0 \Rightarrow_\tau \mu$, a hyper-step
  $\mu \xrightarrow{l} \mu'$, and $\mu' \Rightarrow_\tau \mu_f$.
\end{definition}

\begin{definition}[Weak closure $sys^w$]
  \label{def:weakclosure}
  \lean{PLTS.System.weakClosure}\leanok
  \uses{def:weaktau, def:weakstep}
  The \emph{weak closure} $sys^w$ has the same state space and initial state as
  $sys$, but replaces its steps by weak transitions from a Dirac: for internal $l$,
  $s \xrightarrow{l} \mu$ iff $weakTau\ sys\ \delta_s\ \mu$; for external $l$,
  $s \xrightarrow{l} \mu$ iff $weakStep\ sys\ \delta_s\ l\ \mu$.
\end{definition}

\begin{definition}[Distribution-monad lift $\mathcal{D}(sys)$]
  \label{def:dist}
  \lean{PLTS.System.dist}\leanok
  \uses{def:hyperstep}
  The \emph{distribution-monad lift} $\mathcal{D}(sys)$ has state space
  $\operatorname{PMF} State$, initial state $\delta_{sys.init}$, and steps
  $\mu \xrightarrow{l} \omega$ (with $\omega : \operatorname{PMF}(\operatorname{PMF} State)$)
  exactly when $hyperStep\ sys\ \mu\ l\ (\omega \mathbin{\gg\!=}\mathrm{id})$. It
  packages ``a distribution over states'' as a single state.
\end{definition}

\section{Probabilistic simulations}

A simulation of a \emph{concrete} system $sys_C$ by an \emph{abstract} system
$sys_A$ is carried by a relation $R$ and a step-matching condition, phrased via a
coupling of the resulting distributions.

\begin{definition}[Coupling of distributions]
  \label{def:pmfrel}
  \lean{PLTS.PMFRel}\leanok
  For $R : \alpha \to \beta \to \mathrm{Prop}$, the lifted relation
  $\operatorname{PMFRel} R\ \mu_1\ \mu_2$ holds when there is a joint distribution
  $\omega$ on $\alpha \times \beta$ with marginals $\mu_1, \mu_2$ and support
  contained in $R$.
\end{definition}

\begin{definition}[Strong probabilistic simulation]
  \label{def:strongsim}
  \lean{PLTS.StrongProbabilisticSimulation}\leanok
  \uses{def:system, def:pmfrel}
  $R : State_C \to State_A \to \mathrm{Prop}$ is a \emph{strong} simulation of
  $sys_C$ by $sys_A$ when $R\ sys_C.init\ sys_A.init$, and every concrete step
  $s_C \xrightarrow{l} \mu_C$ from an $R$-related pair $(s_C, s_A)$ is matched by a
  single strong abstract step $s_A \xrightarrow{l} \mu_A$ (same label, no
  stuttering) with $\operatorname{PMFRel} R\ \mu_C\ \mu_A$.
\end{definition}

\begin{definition}[Weak probabilistic simulation]
  \label{def:weaksim}
  \uses{def:weaktau, def:weakstep, def:pmfrel}
  % Paper-level definition; not a standalone Lean declaration (see the NOTE
  % below: the Lean development phrases weak simulation via
  % ProbabilisticForwardSimulation directly).
  As above, but the abstract side may stutter: a concrete step is matched from
  $\delta_{s_A}$ by a $weakTau$ (internal $l$) or a $weakStep$ with the same label
  (external $l$), again with $\operatorname{PMFRel} R\ \mu_C\ \mu_A$.
\end{definition}

\begin{definition}[Probabilistic forward simulation]
  \label{def:forwardsim}
  \lean{PLTS.ProbabilisticForwardSimulation}\leanok
  \uses{def:weaktau, def:weakstep, def:pmfrel}
  Here $R : State_C \to \operatorname{PMF} State_A \to \mathrm{Prop}$ relates each
  concrete state to a \emph{distribution} over abstract states. The initial
  concrete state is related to some distribution supported on $sys_A.init$, and
  from an $R$-related pair $(s_C, \mu_A)$ every concrete step
  $s_C \xrightarrow{l} \mu_C$ is matched by a weak transition
  $\mu_A \Rightarrow (\omega \mathbin{\gg\!=}\mathrm{id})$
  ($weakTau$ for internal $l$, $weakStep$ for external $l$), where
  $\omega : \operatorname{PMF}(\operatorname{PMF} State_A)$ couples $\mu_C$ with
  $R$-related abstract distributions, i.e. $\operatorname{PMFRel} R\ \mu_C\ \omega$.
\end{definition}

Forward simulation generalises weak simulation: it lets the abstract state be
tracked only up to a \emph{belief} distribution, which is precisely what the state
space $\operatorname{PMF} State_A$ of $\mathcal{D}(sys_A^w)$ records.

\section{Building blocks: functional simulation and the constructions}

The soundness proof rests on one reusable engine and three preservation results.

\begin{lemma}[Functional simulation is trace-sound]
  \label{lem:map}
  \lean{PLTS.achievableTraceDists_map}\leanok
  \uses{def:achievable}
  Let $f : X \to Y$ be a label-preserving \emph{functional} simulation between
  systems: $f(sys_X.init) = sys_Y.init$ and every step $s \xrightarrow{l} \mu$
  lifts to $f(s) \xrightarrow{l} f_\ast\mu$ (the pushforward $\mu.map\ f$). Then
  $\operatorname{ATD}(sys_X) \subseteq \operatorname{ATD}(sys_Y)$.
\end{lemma}

\begin{proof}\leanok
  Fix a probabilistic execution $pe_X$ of $sys_X$ with initial distribution
  $\delta_{sys_X.init}$ realising $D$. Since $f$ preserves labels, pushing each
  $X$-history $E_X$ forward to the $Y$-history $f_\ast E_X$ (apply $f$ to every
  state, leave labels alone) preserves the label sequence, hence the trace and the
  tightness status. What must be produced is a \emph{single} $sys_Y$-scheduler
  realising this pushforward, and it is recovered by \emph{disintegration}. To a
  finite $Y$-history $E_Y$ attach the fibre weight
  \[
    g(E_Y, E_X) \;=\; pe_X.probOf(E_X)
    \qquad\text{for } E_X \text{ in the fibre } f_\ast E_X = E_Y,
  \]
  and let the $Y$-scheduler at $E_Y$ emit each next move with probability
  proportional to the total $g$-weight of the fibre members emitting a preimage of
  that move. This is exactly the posterior over the $f$-fibre of $E_Y$ weighted by
  path mass --- it averages $pe_X$'s next-step law over all $X$-histories that $f$
  collapses onto $E_Y$. Because labels are untouched, the induced $sys_Y$-execution
  assigns each tight $Y$-history the summed $probOf$-mass of its fibre, so
  $traceProb$ is preserved cone by cone and the realised trace distribution is
  again $D$.
\end{proof}

\begin{theorem}[Soundness of strong simulation]
  \label{thm:strongsound}
  \lean{PLTS.StrongProbabilisticSimulation.achievableTraceDists_subset}\leanok
  \uses{def:strongsim, lem:map}
  If $R$ is a strong simulation of $sys_C$ by $sys_A$, then
  $\operatorname{ATD}(sys_C) \subseteq \operatorname{ATD}(sys_A)$.
\end{theorem}

\begin{proof}\leanok
  \uses{def:strongsim, lem:map}
  A strong simulation is \emph{lockstep} --- each concrete step is matched by an
  abstract step with the same label, no stuttering --- so the proof factors through
  the \emph{coupled product} $\operatorname{simProd} sys_C\ sys_A\ R$
  (\lean{PLTS.simProd}) over $State_C \times State_A$. Its step relation
  $(s_C, s_A) \xrightarrow{l} \omega$ requires that the two marginals
  $\omega.map\ \pi_1$ and $\omega.map\ \pi_2$ be valid $sys_C$- and $sys_A$-steps
  and that $\operatorname{supp}\omega \subseteq R$; its initial state
  $(sys_C.init, sys_A.init)$ is $R$-related by the simulation's $init$ clause. Two
  inclusions compose.

  \emph{Lift} (\lean{PLTS.simProd_achievableTraceDists_superset}):
  $\operatorname{ATD}(sys_C) \subseteq \operatorname{ATD}(\operatorname{simProd}\,sys_C\,sys_A\,R)$.
  From a $sys_C$-execution $pe_C$ build the joint execution
  \lean{PLTS.simJointExec}: drive $pe_C$ on the concrete component, and whenever it
  emits $(l, \mu_C)$ from a joint history whose current pair $(s_C, s_A)$ is
  $R$-related, apply $sim.step$ to $s_C \xrightarrow{l} \mu_C$ to obtain a coupling
  $\omega$ with $\omega.map\ \pi_1 = \mu_C$, a valid abstract step
  $\omega.map\ \pi_2$, and $\operatorname{supp}\omega \subseteq R$; emit
  $(l, \omega)$. By construction the concrete projection of the joint execution is
  exactly $pe_C$ and labels are shared, so their trace distributions agree. This
  direction is purely forward (no marginalisation).

  \emph{Project}: the second projection $\pi_2$ is a functional, label-preserving
  simulation of $\operatorname{simProd} sys_C\ sys_A\ R$ by $sys_A$ --- it sends the
  joint init to $sys_A.init$, and the marginal condition makes every joint step map
  to a genuine $sys_A$-step $\pi_2(p) \xrightarrow{l} \omega.map\ \pi_2$ --- so
  Lemma~\ref{lem:map} gives
  $\operatorname{ATD}(\operatorname{simProd}\,sys_C\,sys_A\,R) \subseteq \operatorname{ATD}(sys_A)$.
  Compose the two inclusions.
\end{proof}

\begin{theorem}[The weak closure preserves trace distributions]
  \label{thm:weakclosure-eq}
  \lean{PLTS.weakClosure_traceProb_eq}\leanok
  \uses{def:achievable, def:weakclosure}
  $\operatorname{ATD}(sys) = \operatorname{ATD}(sys^w)$.
\end{theorem}

\begin{proof}\leanok
  \uses{lem:map}
  The easy inclusion $\operatorname{ATD}(sys) \subseteq \operatorname{ATD}(sys^w)$
  (\lean{PLTS.weakClosure_traceProb_subset}) is Lemma~\ref{lem:map} with $f = id$:
  every strong $sys$-step is, in particular, a (degenerate) weak step of $sys^w$.

  The hard inclusion (\lean{PLTS.weakClosure_traceProb_superset}) turns an abstract
  scheduler $ws$ on $sys^w$ into a concrete scheduler on $sys$ with the same trace
  distribution. Algorithm~\ref{alg:expand} is the unfolding: it maintains an
  abstract execution $we$ on $sys^w$ and a concrete execution $e$ on $sys$; each
  outer iteration draws a weak transition $wt = (l, \mu)$ from $ws(we)$, takes a
  \emph{witness} scheduler realising it (\lean{PLTS.Realises}: run from
  $\delta_{last(e)}$ it halts almost surely, its halting end-state distribution is
  exactly $\mu$, and every halting trajectory has external trace $[l]$, or $[]$ when
  $l = \tau$), runs that witness to produce the concrete segment
  $e' = \tau^\ast\!\cdot l\cdot\tau^\ast$, and appends it. The control flow is
  deterministic but the draws are random, so each reachable \emph{configuration}
  $c$ is reached with a well-defined probability $reachProb(c)$
  (\lean{PLTS.reachProb}); the induced memoryless concrete scheduler is
  $expandSched\ ws$ (\lean{PLTS.expandSched}), whose next-move law at a concrete
  history is the Bayesian posterior over the next inner draw given that history.

  The fidelity equality \lean{PLTS.traceProb_expandSched_eq} is proved
  configuration by configuration (for $\tau \ne []$; $\tau = []$ is $1 = 1$). Each
  side reduces to a sum of $reachProb$:
  \begin{itemize}
    \item the abstract $traceProb$ collects the \emph{entry} configurations
      $\{\,c \mid c.e'.trans = nil,\ IsTight(c.we),\ trace(c.we) = \tau\,\}$
      (a fresh weak step has just committed);
    \item the concrete $traceProb$ collects the \emph{arrival} configurations
      $\{\,c \mid c.e'.trans \ne nil,\ IsTight(c.e),\ trace(c.e) = \tau\,\}$
      (the concrete run has just emitted the external $l$).
  \end{itemize}
  These are the \emph{same} external weak step measured one moment apart: between an
  arrival and its matching entry lies only the post-$l$ $\tau$-run and the outer
  commit, which change $last(e)$ but not the external trace. Their $reachProb$
  masses are \emph{equal} --- not merely $\le$ --- because the witness's post-$\tau$
  run halts almost surely (the $\sum haltMass = 1$ clause of $Realises$), so no mass
  leaks to divergence inside a witness run and $reachProb$ is conserved across the
  post-$\tau$ sojourn. Hence the entry- and arrival-sums agree.
\end{proof}

\begin{algorithm}
    \caption{Simulating a scheduler on $sys^w$}
    \label{alg:expand}
    \begin{algorithmic}
        \State \textbf{Input:} a PLTS $sys$ and a scheduler $ws$ on $sys^w$
        \medskip
        \State $we, e \gets sys^w.init, sys.init$
        \State $wt \gets ws(we)$
        \State \textbf{while} $wt \ne \bot$:
        \State\quad $s \gets wt.choose$ \Comment{get the scheduler inducing $wt$}
        \State\quad $e' \gets last(e)$
        \State\quad $t \gets s(e')$
        \State\quad \textbf{while} $t \ne \bot$:
        \State\quad\quad $e' \gets e'.lbl(t).out(t)$
        \State\quad\quad $t \gets s(e')$
        \State\quad $e \gets e^\frown e'$
        \State\quad $we \gets we.lbl(wt).last(e)$
        \State\quad $wt \gets ws(we)$
    \end{algorithmic}
\end{algorithm}

\begin{theorem}[The distribution-monad lift preserves trace distributions]
  \label{thm:dist-eq}
  \lean{PLTS.dist_traceProb_eq}\leanok
  \uses{def:achievable, def:dist}
  $\operatorname{ATD}(sys) = \operatorname{ATD}(\mathcal{D}(sys))$.
\end{theorem}

\begin{proof}\leanok
  \uses{lem:map}
  The easy inclusion
  $\operatorname{ATD}(sys) \subseteq \operatorname{ATD}(\mathcal{D}(sys))$
  (\lean{PLTS.dist_traceProb_subset}) is Lemma~\ref{lem:map} for the Dirac
  embedding $s \mapsto \delta_s$: a step $s \xrightarrow{l} \mu$ lifts to the
  hyper-step $\delta_s \xrightarrow{l} \delta_\mu$, whose flattening is $\mu$.

  The hard inclusion (\lean{PLTS.dist_traceProb_superset}) flattens the genuine
  $\operatorname{PMF}(\operatorname{PMF} State)$-randomness of a
  $\mathcal{D}(sys)$-execution $pe'$ back to an ordinary $sys$-execution. Its
  initial distribution is $pe'.initState \mathbin{\gg\!=}\mathrm{id}$ and its
  scheduler is the \emph{lowering} \lean{PLTS.ProbabilisticExecution.lower}: at a
  concrete history $e$ with label list $labs$ and current state $s$, it forms the
  \emph{trace-cone belief} \lean{PLTS.ProbabilisticExecution.beliefTC} --- the
  posterior over those $\mathcal{D}(sys)$-histories $E$ whose full label list is
  $labs$, weighted by the mass that $E$'s current distribution places on $s$ --- and
  emits the corresponding average of $pe'$'s hyper-steps, each split into per-state
  $sys$-steps by \lean{PLTS.ProbabilisticExecution.distHyperKernel}. Conditioning on
  the \emph{full} label list (internal labels included, not just the external trace)
  keeps the belief normalizer finite --- it is a level sum --- and makes the step a
  genuine posterior. The resulting $sys$-execution has the same trace distribution as
  $pe'$. The Dirac-init clause of $\operatorname{ATD}$ is exactly what makes this
  reconstruction sound: a mixture initial distribution over
  $\operatorname{PMF} State$ need not be realised by any single $sys$-execution, and
  the belief reconstruction would then be unsound.
\end{proof}

\section{Soundness of probabilistic forward simulation}

The engine of the reduction is that forward simulation is not a new proof
obligation at all: it is \emph{definitionally} a strong simulation into
$\mathcal{D}(sys_A^w)$. The one non-formal ingredient is the transition-level
correspondence below.

\subsection{The transition-level correspondence}

A step of $\mathcal{D}(sys_A^w)$ from a distribution $\mu$ is a hyper-step of
$sys_A^w$, and it must be identified with a single weak transition of $sys_A$ from
$\mu$. That identification rests on three families of ``workhorse'' lemmas about
how weak transitions interact with mixing (the concrete/abstract subscripts are
dropped; $sys$ is an arbitrary system).

\begin{lemma}[Source decomposition]
  \label{lem:decomp}
  \lean{PLTS.weakTau_exists_pointwise, PLTS.weakStep_exists_pointwise}\leanok
  \uses{def:weaktau, def:weakstep}
  A weak transition out of a distribution factors pointwise: if
  $weakTau\ sys\ \mu\ \nu$ then there is
  $\rho : State \to \operatorname{PMF} State$ with
  $\nu = \mu \mathbin{\gg\!=} \rho$ and $weakTau\ sys\ \delta_s\ (\rho\,s)$ for every
  $s \in \operatorname{supp}\mu$; likewise for $weakStep$.
\end{lemma}

\begin{proof}\leanok
  Dispatch the single witness scheduler on the (observable) start state: run from
  $\delta_s$ it defines $\rho\,s$, and averaging over $\mu$ recovers $\nu$.
\end{proof}

\begin{lemma}[Source mixing]
  \label{lem:mix}
  \lean{PLTS.weakTau_of_pointwise, PLTS.weakStep_of_pointwise}\leanok
  \uses{def:weaktau, def:weakstep}
  The converse recombination: given per-point weak transitions
  $weakTau\ sys\ \delta_s\ (\rho\,s)$ for $s \in \operatorname{supp}\mu$, a single
  scheduler dispatching on the start state witnesses
  $weakTau\ sys\ \mu\ (\mu \mathbin{\gg\!=} \rho)$; likewise for $weakStep$.
\end{lemma}

\begin{lemma}[Target convexity]
  \label{lem:bind}
  \lean{PLTS.weakTau_bind, PLTS.weakStep_bind}\leanok
  \uses{def:weaktau, def:weakstep}
  A \emph{mixture of targets from a single point} is again a weak transition: if
  $q : \operatorname{PMF}(\operatorname{PMF} State)$ satisfies
  $weakTau\ sys\ \delta_s\ \sigma$ for every $\sigma \in \operatorname{supp} q$,
  then $weakTau\ sys\ \delta_s\ (q \mathbin{\gg\!=} \mathrm{id})$; likewise for
  $weakStep$.
\end{lemma}

\begin{proof}\leanok
  The crux. Use a posterior-weighted mixing scheduler: from $\delta_s$ first draw a
  target $\sigma \sim q$ and then run $\sigma$'s witness scheduler. Almost-sure
  halting of each witness gives almost-sure halting of the mixture, and its halting
  end-state distribution is $q \mathbin{\gg\!=} \mathrm{id}$.
\end{proof}

\begin{lemma}[Hyper-step / weak-transition bridge]
  \label{lem:bridge}
  \lean{PLTS.hyperStep_weakClosure_of_weakTau, PLTS.weakTau_of_hyperStep_weakClosure, PLTS.hyperStep_weakClosure_of_weakStep, PLTS.weakStep_of_hyperStep_weakClosure}\leanok
  \uses{def:hyperstep, def:weakclosure, lem:decomp, lem:mix, lem:bind}
  For a distribution $\mu$ over states,
  \[
    hyperStep\ sys^w\ \mu\ l\ \nu
    \iff
    \begin{cases}
      weakTau\ sys\ \mu\ \nu & (l = \tau),\\
      weakStep\ sys\ \mu\ l\ \nu & (l \ne \tau).
    \end{cases}
  \]
  Equivalently: a step of $\mathcal{D}(sys^w)$ out of $\mu$ is exactly a weak
  transition of $sys$ out of $\mu$.
\end{lemma}

\begin{proof}\leanok
  \uses{lem:decomp, lem:mix, lem:bind}
  ($\Leftarrow$) Decompose the weak transition out of $\mu$ into per-point weak
  transitions (Lemma~\ref{lem:decomp}); each is precisely a step of $sys^w$ out of a
  Dirac, so their $\mu$-average is a hyper-step $hyperStep\ sys^w\ \mu\ l\ \nu$.
  ($\Rightarrow$) A hyper-step provides, for each $s \in \operatorname{supp}\mu$, a
  distribution $p\,s$ over $sys^w$-successors of $s$ --- i.e.\ a \emph{mixture of
  weak transitions out of $\delta_s$}. Collapse that mixture to a single per-point
  weak transition with Lemma~\ref{lem:bind}, then recombine over the source with
  Lemma~\ref{lem:mix}.
\end{proof}

\subsection{The reduction}

\begin{theorem}[Forward simulation $=$ strong simulation into $\mathcal{D}(\cdot^w)$]
  \label{thm:forward-equiv}
  \lean{PLTS.probabilisticForwardSimulation_iff_strong_dist_weakClosure}\leanok
  \uses{def:forwardsim, def:strongsim, def:dist, def:weakclosure, lem:bridge}
  For any relation $R : State_C \to \operatorname{PMF} State_A \to \mathrm{Prop}$,
  \[
    \text{$R$ is a forward simulation of $sys_C$ by $sys_A$}
    \iff
    \text{$R$ is a strong simulation of $sys_C$ by $\mathcal{D}(sys_A^w)$.}
  \]
\end{theorem}

\begin{proof}\leanok
  \uses{lem:bridge}
  Both notions carry the \emph{same} relation $R$: the state type of
  $\mathcal{D}(sys_A^w)$ is $\operatorname{PMF} State_A$, exactly the type of the
  forward simulation's abstract argument. On initial states, the forward
  simulation's ``distribution supported on $sys_A.init$'' is forced to be the Dirac
  $\delta_{sys_A.init}$, which is the initial state of $\mathcal{D}(sys_A^w)$; and
  the internal vs.\ external case split coincides because $\tau$ is canonical. The
  step clauses then match verbatim under the bridge (Lemma~\ref{lem:bridge}): a
  forward step matches $s_C \xrightarrow{l} \mu_C$ by a weak transition
  $\mu_A \Rightarrow (\omega \mathbin{\gg\!=} \mathrm{id})$, while a
  $\mathcal{D}(sys_A^w)$ strong step matches it by a hyper-step
  $\mu_A \xrightarrow{l} \omega$, i.e.\
  $hyperStep\ sys_A^w\ \mu_A\ l\ (\omega \mathbin{\gg\!=}\mathrm{id})$ --- and these
  are equivalent, with the \emph{same} coupling $\omega$ and the same
  $\operatorname{PMFRel} R\ \mu_C\ \omega$ witnessing both.
\end{proof}

\begin{theorem}[Soundness of probabilistic forward simulation]
  \label{thm:forward-sound}
  \lean{PLTS.ProbabilisticForwardSimulation.achievableTraceDists_subset}\leanok
  \uses{def:forwardsim, def:achievable, thm:forward-equiv, thm:strongsound,
        thm:dist-eq, thm:weakclosure-eq}
  If $sys_C$ is forward-simulated by $sys_A$ (along some
  $R : State_C \to \operatorname{PMF} State_A \to \mathrm{Prop}$), then
  \[
    \operatorname{ATD}(sys_C) \subseteq \operatorname{ATD}(sys_A).
  \]
\end{theorem}

\begin{proof}\leanok
  \uses{thm:forward-equiv, thm:strongsound, thm:dist-eq, thm:weakclosure-eq}
  Chain the four building blocks. By Theorem~\ref{thm:forward-equiv} the forward
  simulation $R$ is a strong simulation of $sys_C$ by $\mathcal{D}(sys_A^w)$, so
  Theorem~\ref{thm:strongsound} gives
  \[
    \operatorname{ATD}(sys_C) \subseteq \operatorname{ATD}(\mathcal{D}(sys_A^w)).
  \]
  Theorem~\ref{thm:dist-eq}, applied to the system $sys_A^w$, gives
  $\operatorname{ATD}(\mathcal{D}(sys_A^w)) = \operatorname{ATD}(sys_A^w)$, and
  Theorem~\ref{thm:weakclosure-eq}, applied to $sys_A$, gives
  $\operatorname{ATD}(sys_A^w) = \operatorname{ATD}(sys_A)$. Composing,
  \[
    \operatorname{ATD}(sys_C)
      \subseteq \operatorname{ATD}(\mathcal{D}(sys_A^w))
      = \operatorname{ATD}(sys_A^w)
      = \operatorname{ATD}(sys_A). \qedhere
  \]
\end{proof}

The two transformations do all the work: $\cdot^w$ turns the stuttering allowed by
a weak/forward simulation into ordinary strong steps, and $\mathcal{D}(\cdot)$
turns the ``distribution over abstract states'' bookkeeping of a forward
simulation into ordinary states --- and, crucially,
Theorems~\ref{thm:weakclosure-eq} and~\ref{thm:dist-eq} guarantee that neither
transformation changes the observable trace distributions. The analogous soundness
statement for \emph{weak} simulation
% NOTE: PLTS.WeakProbabilisticSimulation is described in SimDefs.lean docstrings but
% not (yet) declared in the code; the \lean{} cross-link is demoted to \texttt{} so
% checkdecls stays green. Restore the \lean{} macro when the declaration lands.
(\texttt{PLTS.WeakProbabilisticSimulation.achievableTraceDists\_subset}) is the same
argument using only $\cdot^w$ (Theorem~\ref{thm:weakclosure-eq}), since a weak
simulation is a strong simulation into $sys_A^w$.


\section{Case study: Asynchronous Byzantine Agreement}

This section tracks the protocol case study built on top of the five results:
the correctness of ABDY22's Asynchronous Byzantine Agreement, with the weak
common coin held at specification level. The development lives in the opt-in
\texttt{Leslie2Protocols} library.

\begin{definition}[Parameters and coin]
  \label{def:aba-params}
  \lean{PLTS.ABA.Params}\lean{PLTS.ABA.Params.coinPMF}\leanok
  The parameters: $n$ processes, at most $f$ corrupted with $3f < n$, and coin
  goodness $\epsilon$ with $2\epsilon \le 1$. One coin resolution follows
  $\mathrm{coinPMF}$: each bit with probability $\epsilon$, the adversarial
  outcome $\top$ otherwise.
\end{definition}

\begin{definition}[Shared alphabet]
  \label{def:aba-labels}
  \lean{PLTS.ABA.Lab}\lean{PLTS.ABA.Lab.hiddenAPI}\leanok
  \uses{def:aba-params}
  All systems of the case study share one label type: the ABA API, the
  round-tagged GBCA/WCC APIs (hidden in the hybrids), corruption
  $\mathrm{fail}(id)$, and $\tau$.
\end{definition}

\begin{definition}[Idle padding and instance families]
  \label{def:idle-family}
  \lean{PLTS.System.withIdle}\lean{PLTS.System.family}\leanok
  Idle self-loops on foreign labels make the full-synchronisation
  \lean{PLTS.System.parallel} emulate sync-set composition; the
  $\mathbb{N}$-indexed instance family moves one instance per owned or silent
  label, broadcasts corruption to all instances, and idles otherwise. Families
  of LTS instances are LTS (\lean{PLTS.System.family_isLTS}\leanok).
\end{definition}

\begin{definition}[The ABA specification, repaired]
  \label{def:aba-spec}
  \lean{PLTS.ABA.SpecStep}\lean{PLTS.ABA.spec}\leanok
  \uses{def:aba-labels}
  Transition System 1 of the ABA blueprint, ten rules, with four deviations.
  Corruption is a total deterministic transform (D1).
  A ghost map $input$ records genuine $\mathrm{callABA}$ events: rule 1 writes it unconditionally, the input-enabledness loop rule 2 writes it first-write-wins (D13).
  The mixed-bind rule 4 carries an $f{+}1$-support guard, requiring $f{+}1$ callers of the bound bit, counted $F$-blind (D13).
  The re-propose rule 7 carries the guard $(val = \bot \wedge (input\,id = b \vee bind = b)) \vee val = b$: while undecided a process re-proposes only a recorded own input or the bound value, and once the decision value is fixed it re-proposes only $val$ (D3, D13).
  A tenth rule, $\mathrm{callByzFill}$, lets a corrupted process fill an empty call slot with an arbitrary bit and no ghost record (D13), the abstract image of the concrete adversary's hidden Byzantine $\mathrm{callG}$ drivers.
  The blueprint's $Initial$ clause mentions an undeclared $out$ field and is omitted.
  The guards are load-bearing.
  Without the rule-7 guard the specification violates Agreement: after a unanimous-$v$ round sets $val := v$ and a process returns $v$, an arbitrary re-proposal of $1-v$ lets the unanimity rule overwrite $val := 1-v$ and a second process return $1-v$.
  Without the rule-4 support guard the specification violates Validity: at $n = 4$, $f = 1$ with inputs $1,0,0,0$ the value $1$ binds off its sole holder, that holder is then corrupted, and $1$ is re-proposed everywhere and decided, yet every never-corrupted process inputs $0$.
  At most $f$ processes are ever corrupted, so $f{+}1$ distinct supporters always include a never-corrupted one, and $3f < n$ supplies $n - 2f \ge f + 1$.
\end{definition}

\blueprintfig{fig-ts1-rules}{Transition System 1, the repaired ABA specification: the three control states, the ten rules as labelled transitions, and the provenance guards on rules 4 and 7.
Rule 4 (\texttt{mixed}) binds under a quorum on both bits with an $f{+}1$-support guard on the bound bit; rule 3 (\texttt{unanim}) decides, either directly from \emph{unbound} or from \emph{bound}, with $val\leftarrow\,!b$ when no honest process holds $b$.
Rule 9 (\texttt{fail}) is a global corruption broadcast to all instances (D1), and rule 10 (\texttt{callByzFill}) lets a corrupted $id\in F$ fill an empty call slot with no ghost record (D13); both are omitted from the diagram.
The provenance guards on rules 4 and 7 exclude the unanimity-overwrite (D3) and later-corrupted-input (D13) traces.}

\begin{theorem}[Safety of the ABA specification]
  \label{thm:aba-spec-safe}
  \lean{PLTS.ABA.spec_safe}\leanok
  \uses{def:aba-spec, thm:trace-support}
  Every trace in the support of every achievable trace distribution of the repaired specification satisfies Validity (\lean{PLTS.ABA.ValidityTrace}\leanok) and Agreement (\lean{PLTS.ABA.AgreementTrace}\leanok).
  Validity is the paper form: every return $\mathrm{retABA}\,id\,b$ at trace position $m$ is preceded by a call $\mathrm{callABA}\,id'\,b$ at some $k < m$ whose caller $id'$ is never corrupted anywhere along the trace.
  Agreement is that any two returns carry the same bit.
  The proof is invariant induction along genuine executions.
  The decision value is write-once: the unanimity rule can only rewrite it to itself, by the quorum argument $f < n - f$, exactly the point where the D3 guard is load-bearing.
  Provenance is a label-history-aware invariant that attributes the decision value to genuine $\mathrm{callABA}$ events: its $f{+}1$ $F$-blind support, minus the at most $f$ ever-corrupted processes, leaves a never-corrupted caller by a budget pigeonhole.
\end{theorem}

\begin{proof}\leanok
  \uses{def:aba-spec, thm:trace-support}
  By Theorem~\ref{thm:trace-support} it suffices to prove the two predicates on the events of every genuine execution, which we do by induction over the ten rules with the invariants $\mathrm{SpecInv}$ and $\mathrm{ValInv}$.
  For Agreement, $\mathrm{SpecInv}$ makes the decision value $val$ write-once: the unanimity rule can rewrite $val$ only to itself, since a $b$-quorum and a $(1-b)$-quorum would meet under $f < n - f$, and the rule-7 (D3) guard blocks the re-proposal that would otherwise flip it.
  For Validity, $\mathrm{ValInv}$ carries provenance conjuncts tying each $bind$ and $val$ to $f{+}1$ $F$-blind recorded genuine inputs of the value.
  The at most $f$ ever-corrupted processes, obtained as a trace-level fold of the $\mathrm{failSet}$, leave one of the $f{+}1$ supporters a never-corrupted recorded caller by a budget pigeonhole.
  Positionality of that caller, a call at some $k < m$, is read off the same label-history fold.
\end{proof}

\begin{definition}[Sub-protocol specifications]
  \label{def:aba-subspecs}
  \lean{PLTS.ABA.WCC.Step}\lean{PLTS.ABA.WCC.specFamily}\lean{PLTS.ABA.GBCA.Step}\lean{PLTS.ABA.GBCA.specFamily}\leanok
  \uses{def:aba-labels, def:idle-family}
  Transition Systems 2 and 3 (round-indexed instances and their families):
  GBCA with linearised binding and graded returns, WCC with the
  $\epsilon$-probabilistic coin resolution. The GBCA side is an LTS
  (\lean{PLTS.ABA.GBCA.specFamily_isLTS}\leanok).
  The GBCA binding and its $B$- and $C$-returns each carry an $f{+}1$ $F$-blind support guard (SuppOK): $f{+}1$ distinct callers-or-corrupted of the relevant bit.
  Since at most $f$ processes are ever corrupted, one of the $f{+}1$ is a never-corrupted genuine caller (D15).
\end{definition}

\blueprintfig{fig-gbca-lifecycle}{The GBCA specification instance (Transition System 2): binding under an $f{+}1$ support guard, the graded returns, and the mutually exclusive A/C grade latch.
Binding (\texttt{bindSet}) fires under a quorum with $f{+}1$ $\mathrm{SuppOK}(v)$ support; \texttt{retB} returns on $f{+}1$ dissent leaving the grade unchanged, \texttt{retA} latches the A-side and \texttt{retC} the C-side.
The A- and C-latches are mutually exclusive, \texttt{retA} needing grade $\in\{\bot,\mathrm A\}$ and \texttt{retC} grade $\in\{\bot,\mathrm C\}$; every $\mathrm{SuppOK}$ count is $f{+}1$ distinct callers-or-corrupted (D15) and $F$-blind.
The \texttt{fail} broadcast is omitted.}

\blueprintfig{fig-wcc}{The WCC specification instance (Transition System 3): the $\epsilon$-probabilistic coin resolution and the return of the common bit or an adversary-chosen bit.
The \texttt{flip} guard requires more than $f$ callers counting corrupted, $|\{id\notin F:\ called\}\cup F|>f$, and is the sole probabilistic transition besides Transition System 1.
On $val{=}1$ and $val{=}0$ every process returns the common bit; on $val{=}\top$ each process returns an adversary-chosen bit.}

\begin{definition}[The GBCA implementation instance]
  \label{def:aba-gbca-impl}
  \lean{PLTS.ABA.GBCA.implInst}\leanok
  \uses{def:aba-labels, def:aba-params}
  ABDY22's GBCA implementation as a per-round PLTS (Algorithm 2).
  Each process runs a four-stage message exchange over an authenticated set-based network.
  On a $\mathrm{callG}$ it multicasts $\langle\mathrm{INPUT},b\rangle$; it relays $\langle\mathrm{INPUT},b\rangle$ once $f{+}1$ receipts of it arrive (the relay gate); it multicasts $\langle\mathrm{ECHO},b\rangle$ on an $n{-}f$ INPUT-$b$ quorum, $\langle\mathrm{VOTE},v\rangle$ on an $n{-}f$ ECHO quorum, and $\langle\mathrm{BIND},v\rangle$ on an $n{-}f$ VOTE quorum, each payload a real bit or $\bot$.
  A process returns grade $A\,b$ on an $n{-}f$ BIND-$b$ quorum, $B\,b$ on an $n{-}f$ any-BIND quorum containing $b$ with $f{+}1$ VOTE-$b$s and $|\mathrm{Valid}|>1$, and $C$ on an $n{-}f$ BIND-$\bot$ quorum with $|\mathrm{Valid}|>1$.
  The network is set-based and authenticated (D5): each sender owns one idempotent message set, the adversary schedules delivery ($\mathrm{deliver}$), and a corrupted $j\in F$ injects an arbitrary message ($\mathrm{byz}$).
  A write-once ghost $bound$ linearises the binding (D6): it is fixed the first time a currently honest process multicasts a real-bit BIND, and all three returns are gated on it; a write-only ghost $grade$ records the returned side (D7).
  Protocol sends are gated on prior participation (D8).
\end{definition}

\blueprintfig{fig-gbca-impl}{The GBCA implementation instance (Algorithm 2): the per-process stage automaton and, boxed, the authenticated set-based network.
On a call a process multicasts $\langle\mathrm{INPUT},b\rangle$ and relays it after $f{+}1$ receipts; an $n{-}f$ quorum at each stage produces the next multicast $\langle\mathrm{ECHO},b\rangle$, $\langle\mathrm{VOTE},v\rangle$, then $\langle\mathrm{BIND},v\rangle$.
The returns are $A\,b$ on an $n{-}f$ BIND-$b$ quorum, $B\,b$ on an $n{-}f$ any-BIND quorum with $f{+}1$ VOTE-$b$ and $|\mathrm{Valid}|>1$, and $C$ on an $n{-}f$ BIND-$\bot$ quorum with $|\mathrm{Valid}|>1$; all three are gated on the ghost bound value.
The dashed box is the D5 network: per-sender sent and received message sets, adversary $\mathrm{deliver}$, and Byzantine injection for $j\in F$, alongside the D6 $bound$ and D7 $grade$ ghosts.}

\begin{theorem}[GBCA instance refinement]
  \label{thm:gbca-inst-refines}
  \lean{PLTS.ABA.GBCA.implRefines}\leanok
  \uses{def:aba-gbca-impl, def:aba-subspecs}
  The GBCA implementation instance is forward-simulated, instance by instance, by the GBCA specification instance.
  The LTS-level forward simulation carries an invariant that establishes the $f{+}1$ support counts the specification's guards demand: the first honest relayer of a bound bit has $f{+}1$ senders of the corresponding INPUT, each a genuine holder or corrupted ($input\_supp$).
  The invariant also carries the quorum-intersection arguments ($f < n - f$) and A/C grade exclusivity.
\end{theorem}

\begin{proof}\leanok
  \uses{def:aba-gbca-impl, def:aba-subspecs, def:forwardsim}
  We exhibit an LTS forward simulation of the implementation instance by the specification instance and check it constructor by constructor.
  The relation carries the invariant $input\_supp$: the first honest relayer of a bound bit $b$ has received $\langle\mathrm{INPUT},b\rangle$ from $f{+}1$ distinct senders, each a genuine holder of $b$ or corrupted.
  It also carries the quorum-intersection facts ($f < n - f$) and A/C grade exclusivity.
  Under the relation the three counts the specification's guards demand --- the $f{+}1$ $\mathrm{SuppOK}$ support of $\mathrm{bindSet}$ and of the $B$- and $C$-returns --- are read off $input\_supp$ and the quorum intersections and transported to the matching abstract step.
  Corruption and the network steps match a spec corruption broadcast or a spec stutter.
\end{proof}

\blueprintfig{fig-gbca-sim}{The GBCA instance forward simulation ($\mathrm{instRel}$): implementation rows above, specification-instance states below, dashed verticals the relation.
The implementation call, the network deliveries (including Byzantine injection), the $bound$-ghost event, and the graded returns align to the specification's $\mathrm{call}$, $\mathrm{bindSet}$, and graded returns, with $F$ aligned on both sides.
Every delivery, relay, echo, and vote row maps to a specification stutter; the implementation call maps to the specification call; the $bound$-ghost event maps to $\mathrm{bindSet}$, whose $f{+}1$ support is supplied by the $input\_supp$ invariant; the graded returns map to the graded specification returns.}

\begin{theorem}[Family congruence and the GBCA family refinement]
  \label{thm:gbca-family-refines}
  \lean{PLTS.ForwardSimulation.family}\lean{PLTS.ABA.GBCA.familyRefines}\leanok
  \uses{thm:gbca-inst-refines, def:idle-family}
  Per-instance forward simulations lift to the $\mathbb{N}$-indexed families
  (pointwise relation; broadcast labels matched by the abstract family's own
  broadcast under a corruption-compatibility lemma), and the lift to a
  probabilistic forward simulation is by the LTS bridge, once, at family
  level.
\end{theorem}

\begin{proof}\leanok
  \uses{thm:gbca-inst-refines, def:idle-family, lem:bridge}
  The per-instance simulations of Theorem~\ref{thm:gbca-inst-refines} lift to the $\mathbb{N}$-indexed families under the pointwise relation.
  Owned and silent labels move the one stepping instance and are matched instance-wise; a corruption broadcast is matched by the abstract family's own broadcast, using a corruption-compatibility lemma that the per-instance relation survives a synchronous $\mathrm{fail}$ at every instance.
  The families are LTS, so the LTS bridge (Lemma~\ref{lem:bridge}) turns the family-level LTS forward simulation into a probabilistic one, applied once at family level.
\end{proof}

\begin{definition}[The ABA core]
  \label{def:aba-core}
  \lean{PLTS.ABA.core}\leanok
  \uses{def:aba-labels}
  The core algorithm as a PLTS (blueprint Algorithm 1).
  Each process runs a five-phase handshake cycle $\mathrm{idle} \to \mathrm{toCallG} \to \mathrm{awaitG} \to \mathrm{toCallW} \to \mathrm{awaitW}$ and back to $\mathrm{toCallG}$ of the next round, carrying an estimate adopted from the graded GBCA outcome or, when the estimate is $\bot$, from the coin.
  The per-round bookkeeping between the coin return and the next GBCA call is fused into the $\mathrm{retW}$ step (D10), which multicasts $\langle\mathrm{DECIDED}, b\rangle$ when the round's grade was $A\,b$.
  A DECIDED gossip layer pools the multicasts in per-process sent and received pools; a corrupted process may pool either bit and so equivocates at this layer (D12$'$).
  A process returns $\mathrm{retABA}\,id\,b$ once $n - f$ distinct senders of $b$ have been delivered to it and $b$ is in its own sent pool.
  Corrupted processes additionally carry Byzantine call- and return-driver rules at each handshake (D11).
\end{definition}

\blueprintfig{fig-core-phases}{The ABA core: the per-process five-phase handshake cycle, the fused DECIDED multicast at the A-grade coin return (D10), and the DECIDED pool layer gating $\mathrm{retABA}$.
The idle self-loop carries \texttt{inputLoop} together with the Byzantine driver rules \texttt{callG}/\texttt{retG}/\texttt{callW}/\texttt{retWByz} for $id\in F$ (any phase, no state change, D11); \texttt{retG} sets $est\leftarrow$ grade, and \texttt{retW} increments the round, adopts the coin when $est=\bot$, and multicasts $\langle\mathrm{DECIDED},b\rangle$ on an A-grade return (D10).
The DECIDED pool layer runs per-round $\tau$-loops: \texttt{deliver} (the adversary moves a sender's pooled $\langle\mathrm{DECIDED},b\rangle$ to a receiver), \texttt{echo} ($f{+}1$ distinct receipts multicast), and \texttt{byzDecided} (D12$'$: $id\in F$ pools either bit).
The exit \texttt{retABA}\,$id\,b$ fires once $n-f$ distinct senders of $b$ have been delivered to $id$ and $b$ is in $id$'s own sent pool.}

\begin{definition}[The hybrids]
  \label{def:aba-hybrids}
  \lean{PLTS.ABA.hybridImpl}\lean{PLTS.ABA.hybridSpec}\leanok
  \uses{def:aba-subspecs, def:aba-labels}
  The composed systems $\mathsf{GBCA} \parallel (\mathsf{Core} \parallel
  \mathsf{WCC.Spec})$ with the sub-protocol API hidden, at implementation and
  specification level for the GBCA slot respectively. GBCA occupies the
  refinable slot of the parallel composition.
\end{definition}

\blueprintfig{fig-hybrid}{The hybrids: the composition $\mathsf{GBCA} \parallel (\mathsf{Core} \parallel \mathsf{WCC.Spec})$ under $\mathrm{abstract}\,\mathrm{hiddenAPI}$.
The GBCA family box is the refinable slot, carrying either the implementation family or the specification family (the substitution); the Core and WCC-specification family boxes are fixed.
One illustrated rendezvous $\mathrm{callG}\,r\,id\,b$: the Core emits it, instance $r$ steps on it, and every other instance idles.
The $\mathrm{abstract}\,\mathrm{hiddenAPI}$ frame turns the GBCA/WCC handshake alphabet into $\tau$; the visible API $\mathrm{callABA}/\mathrm{retABA}/\mathrm{fail}$ leaves the frame.}

\begin{theorem}[Substitution]
  \label{thm:aba-substitution}
  \lean{PLTS.ABA.substSim}\lean{PLTS.ABA.substitution}\leanok
  \uses{def:aba-hybrids, thm:gbca-family-refines}
  Replacing the GBCA implementation family by its specification family under
  the fixed context is a probabilistic forward simulation — the parallel and
  abstraction precongruences applied to the family refinement — and hence a
  trace-distribution inclusion, by soundness.
\end{theorem}

\begin{proof}\leanok
  \uses{def:aba-hybrids, thm:gbca-family-refines, thm:forward-sound}
  Apply the parallel and abstraction precongruences to the family refinement of Theorem~\ref{thm:gbca-family-refines}: parallel composition with the fixed context $\mathsf{Core} \parallel \mathsf{WCC.Spec}$ and then hiding the sub-protocol API each preserve probabilistic forward simulation.
  The result is a probabilistic forward simulation of the implementation-side hybrid by the specification-side hybrid.
  Soundness (Theorem~\ref{thm:forward-sound}) turns it into the trace-distribution inclusion.
\end{proof}

\begin{theorem}[The core simulation]
  \label{thm:aba-coresim}
  \lean{PLTS.ABA.coreSim}\leanok
  \uses{def:aba-hybrids, def:aba-spec}
  The specification-side hybrid is forward-simulated by the repaired ABA specification under the relation $Inv \wedge Abs$.
  The abstraction $Abs$ pins the corrupted set $F$, the returns, and $coin = \bot$ — the twin never fires the coin rule — and is two-phase on the twin's $val$.
  In phase 1, before the first return, the twin is unbound: its abstract calls equal the concrete write-once inputs and its ghost is synced.
  In phase 2, after the first return, the board is clear and $bind = val = some\,v$ with $v$ certified by a concrete A-lock.
  Every hidden row and every concrete coin flip is answered by a stutter under a constant coupling.
  The only burst is $\mathrm{decide\_burst}$ at the first $\mathrm{retABA}$: within one weak step it binds, fills, decides, and returns.
  The concrete invariant carries the cross-round Graded Agreement machinery (downward C-lock propagation, upward bind commitment).
\end{theorem}

\begin{proof}\leanok
  \uses{def:aba-hybrids, def:aba-spec, def:weakstep}
  We check the relation $Inv \wedge Abs$ row by row, one step class at a time, verifying that $Inv$ is preserved at each.
  Every hidden row and the coin flip are answered by an $Abs$-frame stutter: the twin holds still under a constant coupling, and since $Abs$ pins $coin = \bot$ the twin never fires the coin rule, so the probabilistic coin row couples to the stutter.
  The first $\mathrm{retABA}$ is answered by the single $\mathrm{decide\_burst}$ weak step, which within one weak transition binds from the $f{+}1$ DECIDED pools, fills the board, decides, and returns.
  A $\mathrm{fail}$ row is answered by the corruption broadcast.
  The concrete invariant supplies the cross-round Graded Agreement machinery --- downward $C$-lock propagation and upward bind commitment --- that certifies the value bound at the burst.
\end{proof}

\blueprintfig{fig-coresim}{The core forward simulation: the two-phase abstract twin, its stutter answering every concrete row --- including the probabilistic coin flip --- under a constant coupling, and the single $\mathrm{decide\_burst}$ weak step aligned under the first $\mathrm{retABA}$.}

\begin{theorem}[Trace-support bridge]
  \label{thm:trace-support}
  \lean{PLTS.exists_exec_of_traceProb_ne_zero}\lean{PLTS.is_exec_of_probOf_ne_zero}\lean{PLTS.safety_transfer}\leanok
  A trace with positive trace probability under a Dirac-initialised
  probabilistic execution is the trace of a genuine execution of the system,
  with trace membership characterised by the execution's events; safety of
  trace supports transfers along inclusion of achievable trace distributions.
\end{theorem}

\begin{proof}\leanok
  \uses{def:probexec, def:traceprob, def:exec, def:achievable}
  Fix a Dirac-initialised probabilistic execution and a trace $t$ with positive trace probability.
  Positive trace probability forces a positive-mass path, and reading its labels off that path yields a genuine execution of the system whose trace is $t$; trace membership is then characterised by that execution's events.
  For the transfer, suppose $\operatorname{ATD}(sys_C) \subseteq \operatorname{ATD}(sys_A)$ and that every trace in the support of every achievable trace distribution of $sys_A$ is safe.
  Each achievable trace distribution of $sys_C$ is one of $sys_A$, so its support is safe as well.
\end{proof}

\begin{theorem}[Refinement and correctness of ABA]
  \label{thm:aba-main}
  \lean{PLTS.ABA.refines}\lean{PLTS.ABA.main}\leanok
  \uses{thm:aba-substitution, thm:aba-coresim, thm:aba-spec-safe}
  Every trace distribution achievable by the implementation-side hybrid is
  achievable by the ABA specification (soundness applied to the substitution
  and core simulations separately, chained by transitivity of set inclusion —
  the transitivity of probabilistic forward simulation is not used), and
  hence every positive-probability trace of the implementation-side hybrid
  satisfies Validity and Agreement. Machine-checked axiom-clean; the composed
  simulation \lean{PLTS.ABA.simComposed} is additionally provided via
  Theorem 2 and is likewise machine-checked axiom-clean (the transitivity
  of probabilistic forward simulation is now fully proven), documented by a
  pinned axiom guard.
\end{theorem}

\begin{proof}\leanok
  \uses{thm:aba-substitution, thm:aba-coresim, thm:aba-spec-safe, thm:trace-support, thm:forward-sound}
  Soundness (Theorem~\ref{thm:forward-sound}) applied to the substitution simulation (Theorem~\ref{thm:aba-substitution}) and to the core simulation (Theorem~\ref{thm:aba-coresim}) gives two trace-distribution inclusions.
  Chaining them by transitivity of set inclusion yields $\operatorname{ATD}(\mathrm{hybridImpl}) \subseteq \operatorname{ATD}(\mathrm{spec})$, so every trace distribution achievable by the implementation-side hybrid is achievable by the specification.
  By the trace-support bridge (Theorem~\ref{thm:trace-support}), safety of the specification's trace supports (Theorem~\ref{thm:aba-spec-safe}) transfers along this inclusion, so every positive-probability trace of the implementation-side hybrid satisfies Validity and Agreement.
  The composed simulation is obtained separately, by transitivity of probabilistic forward simulation.
\end{proof}

\begin{definition}[Non-vacuity]
  \label{def:aba-examples}
  \lean{PLTS.ABA.P4}\leanok
  \uses{def:aba-hybrids}
  Machine-checked witnesses (at $n = 4$, $f = 1$) that the hybrids execute.
  A connected $21$-step positive-probability run of the specification-side hybrid reaches $\mathrm{retABA}\,0\,\mathrm{true}$: three external inputs, three GBCA handshakes, the quorum-gated binding, three A-returns, three coin calls, the single $\epsilon$-mass coin flip on the agreeing bit, three coin returns each multicasting DECIDED, three deliveries, and the return.
  Every step but the flip is Dirac, and the flip's chosen branch has mass $\epsilon > 0$.
  A synchronised corruption broadcast and the first step of the implementation-side hybrid are witnessed separately.
  Membership of a decision trace in $\mathrm{achievableTraceDists}(\mathrm{hybridImpl}\ P4)$ is not established; only the step-level run is.
\end{definition}
